%% kapa: exam=1 round=28 year=2017 period=1 form=A subject=civil_law subject_ko=민법 seq=36 number=36
%% source: Q-Net 감정평가사 기출문제 (cst00309) artlSeq=5199603 file=extracted/1008/5199603/2017년 제28회 감정평가사 1차시험 문제지/1교시 A형.pdf
%% conversion: kordoc 4.24.51 -> tools/gen1.py (auto); review=proofread-sonnet; flags=none
\begin{question}{36}
유치권에 관한 설명으로 옳은 것은? (다툼이 있으면 판례에 따름)
\begin{choices}
  \item 목적물에 대한 점유를 취득한 후 그 목적물에 관한 채권이 성립한 경우 유치권은 인정되지 않는다.
  \item 유치물이 분할 가능한 경우, 채무자가 피담보채무의 일부를 변제하면 그 범위에서 유치권은 일부 소멸한다.
  \item 유치권자가 유치물을 점유함으로써 유치권을 행사하고 있는 동안에는 피담보채권의 소멸시효는 진행되지 않는다.
  \item 유치권자는 특별한 사정이 없는 한 법원에 청구하지 않고 유치물로 직접 변제에 충당할 수 있다.
  \item 공사업자 乙에게 건축자재를 납품한 甲은 그 매매대금채권에 기하여 건축주 丙의 건물에 대하여 유치권을 행사할 수 없다.
\end{choices}
\end{question}
\answer{5}
